\scanpage{5}\mspage{4}
D'autre part, il y a deux \inserted{\uline{\unsure{autres}}} applications
\begingroup
\makeatletter\tagsleft@true\let\veqno\@@leqno\makeatother
\[
\left\{\begin{aligned}
\text{a)}\quad &\unsure{\mathcal{F}|\mathcal{E}}\longrightarrow\mathcal{P}(\deleted{\unsure{\mathcal{F}|F}}M)&
 F&\longmapsto\operatorname{Omb}_{\unsure{M|\deleted{\mathcal F}M|m}}(F)\\
\text{b)}\quad &\mathcal{F}\hookrightarrow\mathcal{P}\mathcal{P}(\deleted{\unsure{\mathcal{F}|F}}M)&
 \text{et }=\quad\unsure{\mathcal{F}|\widetilde{\mathcal{F}}}&\longrightarrow\operatorname{Fig}(M\unsure{]|)}\subset\mathcal{P}\mathcal{P}\unsure{M|(M)}\\
&& F&\mapsto\deleted{\unsure{O}}\operatorname{Mul\unsure{t|\deleted{\unsure{f}}t}omb}_{M}(F),
\end{aligned}\right.\tag*{(1.\deleted{\unsure{4}}6)}
\]
\endgroup
où $\operatorname{Omb}_{M}(F)$ s'appelle l'\uline{\unsure{ombre|ogombre}} \uline{de la} figure $F$ sur $M$, ou la \uline{grande} \uline{ombre} \uline{de} $\mathcal{F}$. \deleted{\unsure{La}}Une multistrate \inserted{$X$} qui est dans la grande ombre d'une figure $F$ est dite ``\uline{raffiner}'' \uline{la figure} $F$ (c'est cette fois déduit d'une relation d'ordre $\ll$ sur $\mathbb{F}$, le raffinement, par la condition $X\ll F$). Il n'est pas clair que cette application $\mathcal{F}\unsure{\relbar|\to}\mathcal{P}(M)$ soit injective, ie. que $\mathcal{F}$ s'identifie à un ens de parties de $M$, via cette application. Par contre il en sera ainsi de (1.6) b), qui est déduit de a) via \deleted{\unsure{$F$|$\mathcal F$}}la \unsure{formation|fonction} de $\widetilde F$, comme (1.3 b) est déduit de (1.3. a):
\begingroup
\makeatletter\tagsleft@true\let\veqno\@@leqno\makeatother
\[
\operatorname{Multomb}(F)=\{X\in\widetilde F\mid\operatorname{Omb}X\},\tag*{(1.\deleted{\unsure{5|6}}7)\quad\deleted{(4.5)}}
\]
\endgroup
alors que $\operatorname{Omb}(X)$ s'explicite par
\begingroup
\makeatletter\tagsleft@true\let\veqno\@@leqno\makeatother
\[
\unsure{\operatorname{Multomb}|\operatorname{MultOmb}}(X)=\{Y\in\deleted{\unsure{\mathcal{F}|F}}M\mid Y\ll X\},\tag*{(1.\deleted{\unsure{6|7}}8)}
\]
\endgroup
analogue à (1.5), où cette fois on a utilisé la relation \inserted{d'ordre} $\ll$ ald $\leq$ -- \deleted{C'est \unsure{la|le}}
